\section{What beamfs brings}\label{sec:brings}

The measurements answer a narrow question. This section places the
answer among the ways a system can get the right bytes back after a bit
flips between the medium and the filesystem.

\begin{table}[t]
\centering
\caption{Getting correct data back after a flipped bit in file data.
``Measured'' means measured in this report; other rows are as
documented.}
\label{tab:position}
\footnotesize
\begin{tabularx}{\columnwidth}{@{}>{\raggedright\arraybackslash}p{2.0cm}
  >{\raggedright\arraybackslash}X >{\raggedright\arraybackslash}p{1.45cm}@{}}
\toprule
filesystem & on a flipped bit in file data & extra space \\
\midrule
ext4, ext3 & wrong bytes returned, no error (measured) & none \\
btrfs, data \texttt{single} & detected, read refused (measured) &
  checksums \\
btrfs, data \texttt{DUP} or RAID1 & detected, served and repaired from
  the other copy~\cite{btrfsautorepair} & a full copy \\
ZFS, \texttt{copies=2} & detected, repaired from the other
  copy~\cite{zfschecksums} & a full copy~\cite{zfsprops} \\
\bfs capsule & corrected on read up to 8 symbols per codeword, measured
  up to \kernMaxPerDecode{}; refused when the decoder fails; beyond 8, a
  wrong decode is possible without \texttt{DATA\_CSUM} & 288 bytes
  per 4\,096 (7.0\,\%) \\
\bottomrule
\end{tabularx}
\end{table}

\paragraph*{The contribution}
\bfs corrects file data inside the filesystem, on one medium, for 288
bytes in 4\,096: 256 of parity, 16 of descriptor, 16 of generation.
Redundancy also returns correct data (\cref{tab:position}), at the price
of a second copy of every block, and for RAID1 of a second device. The
two protect against different damage: a second copy survives anything
that spares it, a capsule only damage within its correction capacity.
They combine, since a \bfs volume can sit on a mirror. Where a second
copy does not fit, on a single small medium in an embedded, edge or
space system, correction inside the filesystem is what remains; that is
the case \bfs is built for. This report argues that case from the
design, not from a measurement against redundancy, which is planned
(\cref{sec:v4}).

Around that core, \bfs brings: correction on every read, not only once a
checksum has failed; an error instead of bytes the decoder could not
vouch for; a 64-entry correction journal in the superblock recording
block, codeword and symbols corrected; a scrubber that decodes blocks no
one reads and writes corrected ones back; interleaving, so that a burst
in the coded bytes is spread over sixteen codewords
(\cref{sec:capsule}). And it comes with a measurement chain others can
repeat: signed manifests, frozen snapshots, and an audit that found a
defect in its own injector (\cref{sec:injector}).

\paragraph*{What it costs}
288 bytes in 4\,096 of capacity. A Reed-Solomon encode on every write:
v2 measured a median write-and-fsync latency 7.84 times that of ext4
(2.73 times at the 99th percentile), under emulation and on an earlier
layout~\cite{desbrieres2026beamfsv2}; the write path has changed since
and v3 measures no performance. A sequential read fetches most blocks
twice (\cref{sec:readpath}). \moduleLines{} lines of code outside the
kernel tree, with buffer heads still under the metadata.

\paragraph*{Where it is not better today}
Pointers: in the default configuration a flipped indirect pointer that
stays in range is not detected (\cref{sec:indirect-design}), where ext4
with \texttt{metadata\_csum} checksums every extent block~\cite{ext4csum}
and btrfs checksums data and metadata by default~\cite{btrfscsum}.
Inodes: one codeword each, so nine consecutive bytes lose the inode and
the file with it, which interleaving cannot help
(\texttt{burst-resistance.md}). Parity: grouped per codeword, so nine
bytes suffice there too (\cref{sec:capsule}). And the project's own
notes say the default indirect-parity mode should not ship as default
until indirect blocks carry a generation (\cref{sec:limits}, item~\ref{lim:indirect}).
